\documentclass[12pt]{ximera}
\title{Homework 1: Preference Schedules, Plurality, and Borda Count}
\begin{document}
\begin{abstract}
Sections 1.1--1.3: converting a preference schedule between formats, and tabulating an
election by the Plurality method and by the Borda Count method --- which, as the second
problem shows, need not agree.
\end{abstract}
\maketitle

\section*{Sections 1.1--1.3: Ballots, Plurality, and Borda Count}

\begin{problem}
To determine the winner of this year's Hogwarts House Cup, Headmaster Umbridge has
decided to use a preference ballot with her $20$ kittens as voters. Each ballot listed
the candidates Gryffindor, Hufflepuff, Ravenclaw, and Slytherin, and preference was
taken by the order each kitten stepped on the candidate names. These
``printed-names format'' ballots were then organized into the preference schedule
shown in Table 1.

\begin{center}
\begin{tabular}{c|ccccc}
 & 6 & 5 & 5 & 3 & 1 \\\hline
Gryffindor & 3rd & 2nd & 3rd & 2nd & 3rd \\
Hufflepuff & 2nd & 1st & 4th & 1st & 1st \\
Ravenclaw  & 1st & 4th & 1st & 3rd & 4th \\
Slytherin  & 4th & 3rd & 2nd & 4th & 2nd
\end{tabular}

\smallskip
Table 1. Preference schedule, printed-names format.
\end{center}

\begin{prompt}
\textbf{(a)} Rewrite the preference schedule in ``standard format'' --- that is, with
1st, 2nd, 3rd, and 4th listed at the left and candidates written in for each ballot
column. Use G, H, R, and S for the candidates.

\begin{center}
\begin{tabular}{c|ccccc}
Preference & 6 & 5 & 5 & 3 & 1 \\\hline
1st & $\answer{R}$ & $\answer{H}$ & $\answer{R}$ & $\answer{H}$ & $\answer{H}$ \\
2nd & $\answer{H}$ & $\answer{G}$ & $\answer{S}$ & $\answer{G}$ & $\answer{S}$ \\
3rd & $\answer{G}$ & $\answer{S}$ & $\answer{G}$ & $\answer{R}$ & $\answer{G}$ \\
4th & $\answer{S}$ & $\answer{R}$ & $\answer{H}$ & $\answer{S}$ & $\answer{R}$
\end{tabular}
\end{center}

\begin{hint}
Work one column at a time. In the first column, which house was marked 1st? Which was
marked 2nd? Read down the original table looking for the rank you want.
\end{hint}

\begin{feedback}
Each column of the original table lists, for each house, the rank that group of voters
gave it. To convert, read each column and place each house in the row matching its
rank.

For the first column ($6$ voters): Ravenclaw was 1st, Hufflepuff 2nd, Gryffindor 3rd,
Slytherin 4th --- so that column reads R, H, G, S. Repeating for the rest:
\begin{center}
\begin{tabular}{c|ccccc}
Preference & 6 & 5 & 5 & 3 & 1 \\\hline
1st & R & H & R & H & H \\
2nd & H & G & S & G & S \\
3rd & G & S & G & R & G \\
4th & S & R & H & S & R
\end{tabular}
\end{center}
\end{feedback}
\end{prompt}

\begin{prompt}
\textbf{(b)} Using the Plurality method, who wins this year's Hogwarts House Cup?

\begin{multipleChoice}
\choice{Gryffindor}
\choice{Hufflepuff}
\choice[correct]{Ravenclaw}
\choice{Slytherin}
\end{multipleChoice}

How many first-place votes did the winner receive? $\answer{11}$

\begin{feedback}
Plurality counts only first-place votes:
\[
\text{R}: 6+5 = 11, \qquad \text{H}: 5+3+1 = 9, \qquad \text{G}: 0, \qquad \text{S}: 0.
\]
Ravenclaw wins with $11$ of the $20$ first-place votes.

Note that Gryffindor and Slytherin receive \emph{no} first-place votes at all, even
though Gryffindor is ranked 2nd or 3rd on every single ballot. Plurality is blind to
everything below the top line.
\end{feedback}
\end{prompt}
\end{problem}

\begin{problem}
Table 2 gives the preference schedule for an election for Best Pickled Herring of
Trollfjord and surrounding glaciers.

\begin{center}
\begin{tabular}{c|ccccc}
Preference & 27 & 16 & 15 & 13 & 6 \\\hline
1st & A & D & B & B & E \\
2nd & B & E & D & D & C \\
3rd & C & C & C & E & D \\
4th & D & A & E & A & A \\
5th & E & B & A & C & B
\end{tabular}

\smallskip
Table 2. Preference schedule for Best Pickled Herring.
\end{center}

\begin{prompt}
\textbf{(a)} What was the total number of voters?

$\answer{77}$

\begin{feedback}
$27+16+15+13+6 = 77$ voters.
\end{feedback}
\end{prompt}

\begin{prompt}
\textbf{(b)} Using the Plurality method, which pickled herring wins the election?

\begin{multipleChoice}
\choice{A}
\choice[correct]{B}
\choice{C}
\choice{D}
\choice{E}
\end{multipleChoice}

\begin{feedback}
First-place votes: A gets $27$, B gets $15+13=28$, D gets $16$, E gets $6$, and C gets
$0$. B wins with $28$.
\end{feedback}
\end{prompt}

\begin{prompt}
\textbf{(c)} Give the complete ranking of these pickled herrings by Plurality, along
with each one's first-place vote count.

1st place: $\answer{B}$ with $\answer{28}$ votes \qquad
2nd: $\answer{A}$ with $\answer{27}$ \qquad
3rd: $\answer{D}$ with $\answer{16}$

\smallskip
4th: $\answer{E}$ with $\answer{6}$ \qquad
5th: $\answer{C}$ with $\answer{0}$

\begin{feedback}
Sorting the first-place counts:
\[ \text{B }28 > \text{A }27 > \text{D }16 > \text{E }6 > \text{C }0. \]
The counts add to $77$, as they must. \checkmark
\end{feedback}
\end{prompt}

\begin{prompt}
\textbf{(d)} Using the Borda Count method, which pickled herring wins the election?

With $5$ candidates, a 1st-place vote is worth $5$ points, 2nd is worth $4$, and so on
down to $1$ point for 5th.

\begin{multipleChoice}
\choice{A}
\choice[correct]{B}
\choice{C}
\choice{D}
\choice{E}
\end{multipleChoice}

The winner's Borda score is $\answer{270}$.

\begin{hint}
For each candidate, go column by column: multiply the number of voters in the column
by the point value of that candidate's position in it, then add across all five
columns.
\end{hint}

\begin{feedback}
Taking B as the example, column by column:
\[
\begin{array}{lll}
\text{27 voters: B is 2nd} & 4\text{ pts} & 27\cdot 4 = 108\\
\text{16 voters: B is 5th} & 1\text{ pt}  & 16\cdot 1 = 16\\
\text{15 voters: B is 1st} & 5\text{ pts} & 15\cdot 5 = 75\\
\text{13 voters: B is 1st} & 5\text{ pts} & 13\cdot 5 = 65\\
\text{6 voters: B is 5th}  & 1\text{ pt}  & 6\cdot 1 = 6
\end{array}
\]
which totals $108+16+75+65+6 = 270$.

Carrying out the same computation for every candidate gives
\[
\text{B } 270, \quad \text{D } 264, \quad \text{A } 220, \quad \text{C } 211, \quad \text{E } 190.
\]
B wins the Borda count too. As a check, the scores total $1155 = 77\cdot 15$, since
each ballot distributes $5+4+3+2+1=15$ points.
\end{feedback}
\end{prompt}

\begin{prompt}
\textbf{(e)} Give the complete ranking of these pickled herrings by Borda count, with
each score.

1st: $\answer{B}$ with $\answer{270}$ \qquad
2nd: $\answer{D}$ with $\answer{264}$ \qquad
3rd: $\answer{A}$ with $\answer{220}$

\smallskip
4th: $\answer{C}$ with $\answer{211}$ \qquad
5th: $\answer{E}$ with $\answer{190}$

\begin{feedback}
The Borda ranking is
\[ \text{B } 270 > \text{D } 264 > \text{A } 220 > \text{C } 211 > \text{E } 190. \]

Compare this with the Plurality ranking from part (c): B, A, D, E, C. The two methods
agree on the winner but disagree everywhere else. Most strikingly, C finished
\emph{last} under Plurality with zero first-place votes, but comes in 4th under Borda
--- and ahead of E --- because C is ranked 2nd or 3rd on most ballots. A, meanwhile,
drops from 2nd to 3rd: A has many first-place votes but is ranked dead last or nearly
last by everyone else.
\end{feedback}
\end{prompt}
\end{problem}

\end{document}
